Um zu zeigen, dass die gegebene Fläche eine reelle Mannigfaltigkeit ist, definieren wir explizit die Karten und überprüfen die Glattheit der Übergangsfunktionen.

Betrachten wir das Rechteck $R = [0, 2] \times [0, 1]$, das durch die Verklebung der gegenüberliegenden Ränder zu einer Fläche vom Geschlecht 2 wird. Wir definieren lokale Karten, die die Struktur dieser Fläche beschreiben.

\begin{align*}
\text{Karte um die Ecke } (0,0): \quad & \varphi_1: U_1 \to \mathbb{R}^2, \quad \varphi_1(x,y) = (x,y), \\
& \text{wobei } U_1 \text{ eine kleine Umgebung von } (0,0) \text{ in } R \text{ ist.} \\
\text{Karte entlang der Kante } [0,2] \times \{0\}: \quad & \varphi_2: U_2 \to \mathbb{R}^2, \quad \varphi_2(x,y) = (x,y), \\
& \text{wobei } U_2 = [0,2] \times (-\epsilon, \epsilon) \text{ für ein kleines } \epsilon > 0. \\
\text{Karte im Inneren des Rechtecks:} \quad & \varphi_3: U_3 \to \mathbb{R}^2, \quad \varphi_3(x,y) = (x,y), \\
& \text{wobei } U_3 \text{ eine offene Menge im Inneren von } R \text{ ist.}
\end{align*}

Die Übergangsfunktionen zwischen diesen Karten sind wie folgt definiert:

\begin{align*}
\varphi_1 \circ \varphi_2^{-1}: \quad & \text{Identität auf } U_1 \cap U_2, \\
\varphi_2 \circ \varphi_3^{-1}: \quad & \text{Identität auf } U_2 \cap U_3, \\
\varphi_3 \circ \varphi_1^{-1}: \quad & \text{Identität auf } U_3 \cap U_1.
\end{align*}

Da die Übergangsfunktionen Identitäten oder Verschiebungen sind, sind sie glatt. Somit ist die resultierende Fläche eine reelle Mannigfaltigkeit.

Um zu bestimmen, ob diese Atlanten eine komplexe Struktur definieren, müssen die Übergangsfunktionen zwischen den Karten holomorph sein. Da die Übergangsfunktionen in diesem Fall lediglich aus Verschiebungen bestehen, sind sie nicht-holomorph, da sie nicht die Cauchy-Riemann-Gleichungen erfüllen. Somit definieren diese Atlanten keine komplexe Struktur auf der Fläche.

%CHECK: Karten explizit definiert und Übergangsfunktionen als glatt gezeigt. Komplexe Struktur ausgeschlossen durch fehlende Holomorphie.